\documentclass[11pt]{article}

\usepackage[margin=1in]{geometry}
\usepackage{amsmath}
\usepackage{booktabs}
\usepackage{array}

\title{Tcost 2x Leveraged ETF Hedge}
\author{}
\date{}

\begin{document}

\maketitle

\section{Problem Statement}

Estimation of the daily cost of rebalancing a 2x leveraged ETF that tracks an index $I$. The fund must achieve daily 2x exposure using only \textbf{physical equities}, 
\textbf{equity index futures} and \textbf{cash/liquid collateral} ($C$).

\subsection{Informal Overview}

At the close on day $t-1$ the fund holds a known $\mathrm{AUM}$, invested in a
physical stock basket, an index future, and cash/margin, sized to deliver
exactly $2\times$ exposure. Overnight the market moves: each stock drifts by its
own return, the future tracks the index, and a new close is struck on day $t$.

Because the ETF must reset to $2\times$ exposure \emph{every day} while respecting
the exchange-margin and risk-buffer constraints, this price move mechanically
forces a rebalancing trade. The manager must buy (if the index rose) or sell (if
it fell) a basket of stocks in fixed proportions, together with a matching
adjustment to the future. The sizes of these trades are pinned down entirely by
the leverage mechanics and margin rules, and turn out to be linear in the
$\mathrm{AUM}$ and in the index return $r$.

The problem therefore boils down to estimating the market-impact cost of trading
a single correlated basket---the stock constituents plus one future. This is
what the \textbf{Bouchaud--Benzaken cross-impact formula} in
Section~\ref{sec:impact} provides.

\subsection{Notation}

Portfolio variables:

\begin{itemize}
  \item $\mathrm{AUM}_t$: assets under management at the start of day $t$.
  \item $S_t^i$: price of stock $i$ at time $t$ (the symbol $S$ denotes stock
  prices only).
  \item $m$: number of stocks in the replication basket ($i = 1,\ldots,m$).
  \item $w_i$: weight of stock $i$ in the basket, $\sum_{i=1}^{m} w_i = 1$;
  stock $i$ holds dollar value $w_i P_t$.
  \item $P_t$: total dollar value of the physical-equity basket.
  \item $F_t$: futures notional exposure.
  \item $M_t$: exchange margin deposit.
  \item $C_t$: internal liquid cash buffer.
  \item $\mu$: exchange margin requirement ($0 < \mu < 1$).
  \item $V_{\mathrm{max}}$: proxy for the maximum expected daily index drop.
  \item $r$: index return over trading day 1.
  \item $r_i = S_2^i/S_1^i - 1$: return of stock $i$ over trading day 1.
  \item $r_P = \sum_{i} w_i r_i$: physical-basket return; $r_P = r$ under exact
  index replication.
\end{itemize}

Market-impact variables:

\begin{itemize}
  \item $N_P, N_F$: aggregate executed notionals for physical equities and
  futures; stock $i$ receives traded notional $w_i N_P$.
  \item $\sigma_i$: annualized volatility of stock $i$.
  \item $\sigma_F$: annualized volatility of the futures/index exposure.
  \item $\sigma_i^{\mathrm{daily}} = \sigma_i / \sqrt{252}$ and
  $\sigma_F^{\mathrm{daily}} = \sigma_F / \sqrt{252}$: daily volatilities.
  \item $T_i, T_F$: average daily turnover for stock $i$ and futures.
  \item $\mathbf{q}$: signed rebalancing-notional vector over all traded assets.
  \item $\mathbf{v}$: average-daily-turnover vector matching $\mathbf{q}$.
  \item $\mathbf{R}$: correlation matrix of all traded assets (renamed from the
  cross-impact $\mathbf{C}$ matrix to avoid collision with the cash variable
  $C_t$).
  \item $\gamma$: market-impact calibration constant.
  \item $\alpha$: impact exponent, typically near $1/2$.
\end{itemize}

\subsection{Constraints}

\begin{center}
\begin{tabular}{@{}ll@{}}
\toprule
Constraint & Formula \\
\midrule
Total exposure & $P_t + F_t = 2 \cdot \mathrm{AUM}_t$ \\
Balance sheet & $P_t + M_t + C_t = \mathrm{AUM}_t$ \\
Exchange margin & $M_t = \mu \cdot F_t$ \\
Risk buffer & $C_t = V_{\mathrm{max}} \cdot F_t$ \\
\bottomrule
\end{tabular}
\end{center}

\section{Initial Portfolio Allocation (Day 1)}

Define $K$ as the required futures exposure per dollar of AUM:
\[
  K = \frac{1}{(1 - \mu) - V_{\mathrm{max}}}
\]

Substituting the margin and buffer rules into the balance sheet gives:
\[
  P_1 + \mu F_1 + V_{\mathrm{max}}F_1 = \mathrm{AUM}_1
  \implies P_1 + (\mu + V_{\mathrm{max}})F_1 = \mathrm{AUM}_1
\]

Using $P_1 = 2\mathrm{AUM}_1 - F_1$, the Day 1 target holdings are:

\begin{center}
\begin{tabular}{@{}ll@{}}
\toprule
Holding & Day 1 value \\
\midrule
Futures notional & $F_1 = K \cdot \mathrm{AUM}_1$ \\
Physical equities & $P_1 = (2 - K) \cdot \mathrm{AUM}_1$ \\
Internal cash buffer & $C_1 = V_{\mathrm{max}} \cdot K \cdot \mathrm{AUM}_1$ \\
Exchange margin deposit & $M_1 = \mu \cdot K \cdot \mathrm{AUM}_1$ \\
\bottomrule
\end{tabular}
\end{center}

\subsection{Case $V_{\mathrm{max}} = \mu$}

If the internal risk buffer equals the exchange margin rate, then
\[
  K = \frac{1}{1 - 2\mu} \quad \text{with } \mu < \frac{1}{2}.
\]

The allocation simplifies to:

\begin{center}
\begin{tabular}{@{}ll@{}}
\toprule
Holding & Value when $V_{\mathrm{max}} = \mu$ \\
\midrule
Futures notional & $F_1 = \frac{\mathrm{AUM}_1}{1 - 2\mu}$ \\
Physical equities & $P_1 = \left(2 - \frac{1}{1 - 2\mu}\right)\mathrm{AUM}_1$ \\
Internal cash buffer & $C_1 = \frac{\mu}{1 - 2\mu}\mathrm{AUM}_1$ \\
Exchange margin deposit & $M_1 = \frac{\mu}{1 - 2\mu}\mathrm{AUM}_1$ \\
\bottomrule
\end{tabular}
\end{center}

\section{Daily Rebalancing}

If the index returns $r$, the 2x fund value becomes:
\[
  \mathrm{AUM}_2 = \mathrm{AUM}_1(1 + 2r)
\]

\subsection{Before Rebalancing}

Before the close, each holding moves by its own return:
\begin{itemize}
  \item Stock $i$: its dollar holding $w_i P_1$ grows to $w_i P_1(1 + r_i)$.
  \item Physical basket (organic): summing over stocks,
  \[
    P_{\mathrm{organic}} = \sum_i w_i P_1(1 + r_i) = P_1(1 + r_P),
  \]
  which equals $P_1(1 + r)$ under exact index replication ($r_P = r$).
  \item Futures: the notional exposure scales to $F_1(1 + r)$, with realized
  cash P\&L $F_1 \cdot r$.
\end{itemize}

\subsection{Target Day 2 Values}

To keep 2x leverage on the new AUM:
\begin{itemize}
  \item $F_2 = K \cdot \mathrm{AUM}_2 = K \cdot \mathrm{AUM}_1(1 + 2r)$.
  \item $P_2 = (2 - K) \cdot \mathrm{AUM}_2
    = (2 - K) \cdot \mathrm{AUM}_1(1 + 2r)$.
\end{itemize}

\subsection{Active Rebalancing Trades}

The closing trades bridge organic end-of-day values and Day 2 targets:

\begin{center}
\renewcommand{\arraystretch}{1.35}
\begin{tabular}{@{}p{0.20\textwidth}p{0.18\textwidth}p{0.18\textwidth}p{0.32\textwidth}@{}}
\toprule
Trade & Organic value & Target value & Executed amount \\
\midrule
Physical equities &
$P_1(1 + r)$ &
$P_2 = P_1(1 + 2r)$ &
$\mathrm{d}P = P_2 - P_1(1 + r) = P_1r
= \left(2 - \frac{1}{(1 - \mu) - V_{\mathrm{max}}}\right)\mathrm{AUM}_1r$ \\
Futures &
$F_1(1 + r)$ &
$F_2 = F_1(1 + 2r)$ &
$\mathrm{d}F = F_2 - F_1(1 + r) = F_1r
= \left(\frac{\mathrm{AUM}_1}{(1 - \mu) - V_{\mathrm{max}}}\right)r$ \\
\bottomrule
\end{tabular}
\end{center}

\section{Market-Impact Transaction Costs}
\label{sec:impact}

Treat the stock constituents and futures contract as one correlated traded asset
universe. Let the signed rebalancing vector be
\[
  \mathbf{q}
  = (q_1,\ldots,q_m,q_F)^\top,
  \qquad
  q_i = w_i\,\mathrm{d}P,
  \qquad
  q_F = \mathrm{d}F.
\]

The per-stock split $q_i = w_i\,\mathrm{d}P$ uses day-1 weights; the intraday
weight drift is $O(r)$, so it only perturbs $q_i$ at order $O(r^2)$ and is
neglected.

The corresponding turnover vector and daily-volatility diagonal matrix are
\[
  \mathbf{v} = (T_1,\ldots,T_m,T_F)^\top,
  \qquad
  \mathbf{D}_{\sigma}
  = \operatorname{diag}(\sigma_1^{\mathrm{daily}},\ldots,
  \sigma_m^{\mathrm{daily}},\sigma_F^{\mathrm{daily}}).
\]

Using the Bouchaud--Benzaquen cross-impact formula, the market-impact cost in bp
for each traded asset is
\[
  \mathbf{I}^{\mathrm{bp}}
  = 10^4\gamma\,\mathbf{D}_{\sigma}\,\mathbf{R}^{1-\alpha}
  \left[
  \operatorname{sign}(\mathbf{q}) \odot
  \left|\mathbf{R}^{\alpha}\left(\frac{\mathbf{q}}{\mathbf{v}}\right)\right|^{\alpha}
  \right].
\]

Division, absolute value, powers, signs, and $\odot$ are element-wise except for
the matrix powers of $\mathbf{R}$. The entries of $\mathbf{q}$ are:

\begin{center}
\small
\renewcommand{\arraystretch}{2.25}
\begin{tabular}{@{}p{0.18\textwidth}p{0.36\textwidth}p{0.40\textwidth}@{}}
\toprule
Rebalanced leg & Signed vector entry & Absolute notional \\
\midrule
Stock $i$ &
$q_i = w_i\left(2 - \frac{1}{(1 - \mu) - V_{\mathrm{max}}}\right)
\mathrm{AUM}_1r$ &
$\lvert q_i\rvert = w_iN_P$ \\
Futures &
$q_F = \frac{\mathrm{AUM}_1r}{(1 - \mu) - V_{\mathrm{max}}}$ &
$N_F = \lvert q_F\rvert$ \\
\bottomrule
\end{tabular}
\end{center}

With this sign convention, the total market-impact cost in USD is
\[
  \operatorname{tcost}^{\mathrm{USD}}
  = \frac{\mathbf{q}^{\top}\mathbf{I}^{\mathrm{bp}}}{10^4}.
\]

Equivalently, by components,
\[
  \operatorname{tcost}^{\mathrm{USD}}
  = \frac{1}{10^4}
  \left(\sum_i q_i I_i^{\mathrm{bp}} + q_F I_F^{\mathrm{bp}}\right).
\]

\section{Operational Conclusion}

Both active trades scale \textbf{linearly} with $\mathrm{AUM}_1$ and $r$.

\begin{itemize}
  \item If $r > 0$, the manager \textbf{buys} both physical equities and futures.
  \item If $r < 0$, the manager \textbf{sells} both physical equities and futures.
\end{itemize}

\end{document}